\documentclass[11pt]{article}
\usepackage[margin=1in]{geometry}
\usepackage{amsmath,amssymb,amsthm,mathtools}
\usepackage{booktabs,graphicx,microtype,enumitem}
\usepackage[hidelinks]{hyperref}

\newtheorem{theorem}{Theorem}[section]
\newtheorem{proposition}[theorem]{Proposition}
\newtheorem{corollary}[theorem]{Corollary}
\newtheorem{lemma}[theorem]{Lemma}
\theoremstyle{remark}
\newtheorem{remark}[theorem]{Remark}

\newcommand{\Q}{\mathbb Q}
\newcommand{\Z}{\mathbb Z}
\newcommand{\F}{\mathbb F}
\newcommand{\OO}{\mathcal O}
\newcommand{\tr}{\operatorname{tr}}
\newcommand{\Gal}{\operatorname{Gal}}
\newcommand{\disc}{\operatorname{disc}}
\newcommand{\Fr}{\operatorname{Fr}}
\newcommand{\Tr}{\operatorname{Tr}}
\newcommand{\Art}{\operatorname{Art}}

\title{\textbf{The Discriminant-$12$ Return}\\
\large Centered Cayley Reciprocity, Pell--Lorentz Geometry, Cyclotomic Closure, and Finite Reduction}
\author{Jeremy Ebert}
\date{2026}

\begin{document}
\maketitle

\begin{abstract}
For $g\in SL_2$ with trace $\tau$, centered part
$H=g-\tfrac{\tau}{2}I$, and Cayley transform
$K=(g-I)(g+I)^{-1}$, we prove
\[
K=\frac{2}{\tau+2}H,
\qquad
KH=\frac{\tau-2}{2}I.
\]
Thus, in the normalization $\tau=\sigma+2$, the hyperbolic self-reciprocity
condition $K=H^{-1}$ selects $\sigma=2$, hence trace $4$.  With the
period-$[1,2]$ positive integral normalization this yields
\[
g_{12}=\begin{pmatrix}3&1\\2&1\end{pmatrix},
\qquad
\lambda=2+\sqrt3,
\]
with characteristic discriminant $12$.  We realize $g_{12}$ as multiplication
by $\lambda$ on the ramified ideal above $2$ in $\Q(\sqrt3)$, and its centered
generator as multiplication by $\sqrt3$.  The same unit acts on factor-cone
null coordinates by $(x,y)\mapsto(\lambda x,\lambda^{-1}y)$ and translates
rapidity by $\log\lambda$ on every real shell $xy=n$.  The row/column
parabola--circle tangent endpoints are its null eigenrays.  After adjoining
$i$, the operator $\mathcal Z=(H_{12}+iI)/2$ is multiplication by
$\zeta_{12}$ on the integral cyclotomic ideal
$\mathfrak p_2\OO_{\Q(\zeta_{12})}$.  At the ramified prime the cyclotomic
residue field is $\F_4$: reduced multiplication by $\zeta_{12}$ has order
$3$, while the reduced Pell return is Frobenius.  Under a compatible additive
identification, the reduced Lorentz norm is exactly the finite-field trace,
so Boolean XOR parity is $\Tr_{\F_4/\F_2}$.  Finally $n=11$ appears as a
distinguished specialization whose Artin action is $T_{11}$, complex
conjugation, corresponding to the involution $J$ in the companion
representation.  No statement about zeta zeros or the Riemann hypothesis is
used.
\end{abstract}

\section{Introduction}

The construction begins before any arithmetic normalization.  Centered Cayley
reciprocity selects trace $4$.  A continued-fraction normalization then
chooses a distinguished primitive positive integral representative.  Its
expanding eigenvalue is the fundamental Pell unit $2+\sqrt3$.  The same
operator is subsequently realized on an ideal lattice, as a Lorentz boost on
the real factor cone, as a cyclotomic operator after adjoining $i$, and as
order-$2$ and order-$3$ finite actions after ramified reduction.

The logical chain is
\[
\boxed{
\begin{gathered}
K=H^{-1}\Longleftrightarrow\sigma=2
\Longrightarrow\tau=4\Longrightarrow\Delta=12\\
\Longrightarrow\lambda=2+\sqrt3
\Longrightarrow g_{12}=\times\lambda\text{ on }\mathfrak p_2
\Longrightarrow H_{12}=\times\sqrt3\\
\Longrightarrow \mathcal Z=\times\zeta_{12}\text{ on }\mathfrak P_2
\Longrightarrow \F_4\text{ reduction}.
\end{gathered}}
\]
The real Cone is a parallel carrier of the same Pell return, not a finite
quotient.  Likewise the Boolean translation group, factor exchange,
$U(12)\cong V_4$, and the affine/semilinear groups of $\F_4$ are distinct
symmetry layers.

\section{Main theorem}

\begin{theorem}[Discriminant-$12$ return package]
\label{thm:main}
Let $g\in SL_2(\Q)$ have trace $\tau$, assume $g+I$ is invertible, and set
\[
H=g-\frac\tau2I,
\qquad K=(g-I)(g+I)^{-1},
\qquad \tau=\sigma+2.
\]
Then
\[
K=\frac{2}{\tau+2}H,
\qquad KH=\frac{\tau-2}{2}I.
\]
If $g$ is hyperbolic, then
\[
K=H^{-1}\iff\sigma=2.
\]
At this selected trace $\tau=4$, impose the period-$[1,2]$ positive integral
normalization and define
\[
g_{12}=\begin{pmatrix}3&1\\2&1\end{pmatrix},
\qquad
H_{12}=g_{12}-2I.
\]
Then the following hold.
\begin{enumerate}[label=(\roman*)]
\item $\chi_{g_{12}}(X)=X^2-4X+1$, $\Delta_{g_{12}}=12$, and
$H_{12}^2=3I$.
\item On $\mathfrak p_2=(2,\sqrt3-1)=(1+\sqrt3)$, $g_{12}$ and $H_{12}$
are multiplication by $2+\sqrt3$ and $\sqrt3$ respectively.
\item The primitive form $q_{12}(x,y)=2x^2-2xy-y^2$ is preserved by
$g_{12}$ and satisfies
\[
2q_{12}=U^2-V^2,
\qquad (U,V)=(2x-y,\sqrt3\,y).
\]
\item In factor-cone null coordinates the same return is
\[
x'=(2+\sqrt3)x,
\qquad y'=(2-\sqrt3)y,
\]
so $xy$ is invariant and rapidity is translated by $\log(2+\sqrt3)$.
The factor-parabola tangent endpoints are the two null eigenrays.
\item If $K_{12}=\Q(\zeta_{12})$ and
$\mathfrak P_2=\mathfrak p_2\OO_{K_{12}}$, then
\[
\mathcal Z=\frac{H_{12}+iI}{2}=\times\zeta_{12}
\]
on $\mathfrak P_2$, and
$[\OO_{K_{12}}:\Z[\sqrt3,i]]=4$.
\item $\OO_{K_{12}}/\mathfrak P_2\cong\F_4$.  Reduced $\mathcal Z$ is
multiplication by an element $\omega$ of order $3$, while reduced $g_{12}$ is
Frobenius.  These generate $GL_2(\F_2)\cong S_3$.
\end{enumerate}
\end{theorem}

\section{Centered Cayley reciprocity}

\begin{lemma}
For $g\in SL_2$ with trace $\tau$,
\[
H^2=\left(\frac{\tau^2}{4}-1\right)I.
\]
\end{lemma}
\begin{proof}
Cayley--Hamilton gives $g^2-\tau g+I=0$.  Substituting
$H=g-\tau I/2$ gives the result.
\end{proof}

\begin{proposition}
If $g+I$ is invertible, then
\[
K=\frac{2}{\tau+2}H,
\qquad KH=\frac{\tau-2}{2}I.
\]
\end{proposition}
\begin{proof}
Cayley--Hamilton gives
\[
(g+I)((\tau+1)I-g)=(\tau+2)I,
\]
whence
\[
(g+I)^{-1}=\frac{(\tau+1)I-g}{\tau+2}.
\]
Therefore
\[
K=\frac{2g-\tau I}{\tau+2}=\frac{2}{\tau+2}H,
\]
and multiplication by $H$ gives the second identity.
\end{proof}

\begin{corollary}
If $g$ is hyperbolic and $\tau=\sigma+2$, then
\[
K=H^{-1}\iff\sigma=2.
\]
\end{corollary}
\begin{proof}
Hyperbolicity makes $H$ invertible, so $K=H^{-1}$ iff $KH=I$, which is
$(\tau-2)/2=1$.
\end{proof}

\section{The primitive integral return and the ramified ideal}

Let
\[
A_a=\begin{pmatrix}a&1\\1&0\end{pmatrix}.
\]
The period $[1,2]$ gives
\[
g_{12}=A_1A_2=\begin{pmatrix}3&1\\2&1\end{pmatrix}.
\]
The Cayley condition selects the trace-$4$ hyperbolic class; the
period/positivity normalization selects this representative.  Directly,
\[
\det g_{12}=1,
\quad \tr g_{12}=4,
\quad \chi_{g_{12}}(X)=X^2-4X+1,
\]
so the eigenvalues are $2\pm\sqrt3$ and the characteristic discriminant is
$12$.  For
\[
H_{12}=g_{12}-2I=\begin{pmatrix}1&1\\2&-1\end{pmatrix},
\]
one has $H_{12}^2=3I$.

Let
\[
f_1=2,
\qquad f_2=\sqrt3-1,
\qquad \mathfrak p_2=(f_1,f_2).
\]
Then
\[
\sqrt3 f_1=f_1+2f_2,
\qquad
\sqrt3 f_2=f_1-f_2,
\]
so
\[
[\times\sqrt3]_f=H_{12},
\qquad
[\times(2+\sqrt3)]_f=g_{12}.
\]

A second natural basis is
\[
e_1=1+\sqrt3,
\qquad e_2=3+\sqrt3=\sqrt3 e_1.
\]
Since $e_1=f_1+f_2$ and $e_2=2f_1+f_2$,
\[
C_{\rm basis}=\begin{pmatrix}1&2\\1&1\end{pmatrix},
\qquad
C_{\rm basis}^{-1}g_{12}C_{\rm basis}
=\begin{pmatrix}2&3\\1&2\end{pmatrix}.
\]
Modulo $2$, this integral basis change will become the Boolean relabeling.

\section{Pell--Lorentz and factor-cone realizations}

For
\[
\alpha=xf_1+yf_2=(2x-y)+y\sqrt3,
\]
\[
N_{F/\Q}(\alpha)
=(2x-y)^2-3y^2
=2q_{12}(x,y),
\]
where
\[
q_{12}(x,y)=2x^2-2xy-y^2.
\]
Thus, with
\[
U=2x-y,
\qquad V=\sqrt3\,y,
\]
\[
2q_{12}=U^2-V^2.
\]
Multiplication by $2+\sqrt3$ is the Lorentz boost
\[
\begin{pmatrix}U'\\V'\end{pmatrix}
=
\begin{pmatrix}2&\sqrt3\\\sqrt3&2\end{pmatrix}
\begin{pmatrix}U\\V\end{pmatrix}.
\]

Now introduce Paper-A factor-cone coordinates
\[
X=\frac{x-y}{2},
\qquad Y=\sqrt{xy},
\qquad T=\frac{x+y}{2}.
\]
Then
\[
T^2-X^2-Y^2=0,
\qquad x=T+X,
\qquad y=T-X.
\]
Consequently
\[
\boxed{x'=\lambda x,\qquad y'=\lambda^{-1}y,}
\qquad \lambda=2+\sqrt3,
\]
and $xy=n$ is invariant.  If
\[
x=\sqrt n e^s,
\qquad y=\sqrt n e^{-s},
\]
then
\[
\boxed{s'=s+\log(2+\sqrt3).}
\]
Factor exchange sends $s\mapsto-s$ and reverses the boost.

For fixed factor level $u>0$, the row parabola is
\[
Y^2=u^2+2uX
\]
with vertex $(-u/2,0)$, while its reflected column parabola has vertex
$(u/2,0)$.  A fixed-$T$ Cone slice is
\[
X^2+Y^2=T^2.
\]
Tangency at the vertices occurs exactly for
\[
T=u/2.
\]
Hence the side-view tangent points are
\[
(T,X)=\left(\frac u2,\pm\frac u2\right),
\]
which are precisely the two null eigenrays of the discriminant-$12$ boost.

\section{The mod-$12$ unit shell}

The units modulo $12$ are
\[
U(12)=\{1,5,7,11\}\cong V_4.
\]
For $r\in U(12)$, the factor pairs $(r,1)$ and $(1,r)$ have
\[
T_r=\frac{r+1}{2},
\qquad X_r=\pm\frac{r-1}{2},
\qquad Y_r=\sqrt r.
\]
Thus the four labels occupy $T=1,3,4,6$, and the outer $r=11$ pair lies on
$x+y=12$.  The standard cyclotomic action is
\[
\Gal(\Q(\zeta_{12})/\Q)\cong U(12),
\qquad
T_r:\zeta_{12}\mapsto\zeta_{12}^{r}.
\]
The geometry carries the labels but does not define their modular group law.

\begin{figure}[ht]
\centering
\includegraphics[width=\textwidth]{mod12_v4_cone_triple.png}
\caption{The mod-$12$ unit labels in factor-cone coordinates.  The outer
$r=11$ pair is $(11,1)$ and $(1,11)$ on $x+y=12$.}
\label{fig:mod12-unit-shell}
\end{figure}

\section{Cyclotomic completion on an integral ideal}

Let
\[
K_{12}=\Q(\zeta_{12})=\Q(\sqrt3,i),
\qquad
\zeta=\frac{\sqrt3+i}{2}.
\]
Since $H_{12}=\times\sqrt3$,
\[
\boxed{\mathcal Z=\frac{H_{12}+iI}{2}=\times\zeta.}
\]
The correct integral carrier is
\[
\mathfrak P_2=\mathfrak p_2\OO_{K_{12}}.
\]
For $\alpha=1+\sqrt3$,
\[
\{\alpha,\alpha\zeta,\alpha\zeta^2,\alpha\zeta^3\}
\]
is a $\Z$-basis of $\mathfrak P_2$.  Since
\[
\Phi_{12}(X)=X^4-X^2+1,
\qquad \zeta^4=\zeta^2-1,
\]
multiplication by $\zeta$ has matrix
\[
\boxed{
\begin{pmatrix}
0&0&0&-1\\
1&0&0&0\\
0&1&0&1\\
0&0&1&0
\end{pmatrix}.}
\]
Thus the half in $(H_{12}+iI)/2$ is a coordinate effect, not a failure of
integrality on the correct lattice.

The full ring of integers is $\OO_{K_{12}}=\Z[\zeta_{12}]$, whereas
$\OO_0=\Z[\sqrt3,i]$ is a suborder of index
\[
\boxed{[\OO_{K_{12}}:\OO_0]=4.}
\]
Equivalently,
\[
\disc(\OO_0)=2304,
\qquad \disc(\OO_{K_{12}})=144.
\]
The additive quotient $\OO_{K_{12}}/\OO_0\cong(\F_2)^2$ records
integral-closure gluing; it is not the residue field below.

\section{Ramified reduction, Frobenius, and XOR trace}

Modulo $2$ in the ramified basis,
\[
\bar g_{12}
=\begin{pmatrix}1&1\\0&1\end{pmatrix},
\qquad
\bar q_{12}(x,y)=y.
\]
Also
\[
C_{\rm basis}\equiv
\Phi=\begin{pmatrix}1&0\\1&1\end{pmatrix}\pmod2,
\]
and therefore
\[
\Phi\bar g_{12}\Phi^{-1}
=P=\begin{pmatrix}0&1\\1&0\end{pmatrix}.
\]
Thus the Boolean relabeling is the reduction of an integral ideal basis
change.  In Boolean coordinates $(\epsilon_1,\epsilon_2)$,
\[
\bar q_{12}=\epsilon_1\oplus\epsilon_2.
\]
Together with Boolean translations, factor exchange gives
$V_4^{\rm tr}\rtimes C_2\cong D_8$.

The cyclotomic residue field is a different four-state ring:
\[
\boxed{\OO_{K_{12}}/\mathfrak P_2\cong\F_4.}
\]
Let $\omega=\bar\zeta$.  Then
\[
\omega^2+\omega+1=0,
\qquad \omega^3=1.
\]
In the additive basis $(1,\omega)$,
\[
[\times\omega]=\begin{pmatrix}0&1\\1&1\end{pmatrix},
\qquad
[\Fr]=\begin{pmatrix}1&1\\0&1\end{pmatrix}.
\]
Hence
\[
\boxed{[\Fr]=\bar g_{12}},
\qquad
\Fr(\times\omega)\Fr^{-1}=(\times\omega)^{-1},
\]
so
\[
\langle\times\omega,\Fr\rangle
=GL_2(\F_2)\cong S_3.
\]
With translations,
\[
\operatorname{AGL}(1,4)\cong A_4,
\qquad
\operatorname{A\Gamma L}(1,4)\cong S_4.
\]

\subsection{Compatible additive identification and finite-field trace}

There is a particularly useful $\F_2$-linear identification between the
ramified four-state additive group and $\F_4$.  Send
\[
\boxed{f_1\longmapsto1,
\qquad f_2\longmapsto\omega^2.}
\]
Because $e_1=f_1+f_2$ and $e_2=f_2$ modulo $2$, this gives
\[
e_1\longmapsto1+\omega^2=\omega,
\qquad
e_2\longmapsto\omega^2.
\]
Thus the principal/Pell basis corresponds to $(\omega,\omega^2)$, a basis
interchanged by Frobenius.  Consequently the same involution is represented
as the transvection $\bar g_{12}$ in the ramified basis and as factor exchange
$P$ in the principal/Boolean basis.

The field trace satisfies
\[
\Tr_{\F_4/\F_2}(1)=0,
\qquad
\Tr_{\F_4/\F_2}(\omega)
=\Tr_{\F_4/\F_2}(\omega^2)=1.
\]
Hence in the ramified basis
\[
\Tr_{\F_4/\F_2}(x+y\omega^2)=y=\bar q_{12}(x,y),
\]
while in the Boolean/Pell basis
\[
\boxed{
\Tr_{\F_4/\F_2}
(\epsilon_1\omega+\epsilon_2\omega^2)
=\epsilon_1\oplus\epsilon_2.}
\]
Therefore, under this compatible additive identification,
\[
\boxed{\bar q_{12}=\Tr_{\F_4/\F_2},}
\]
so the Boolean XOR parity is exactly the finite-field trace.

\begin{remark}[Ring guardrail]
This is an additive $\F_2$-linear identification, not a ring isomorphism.
The ramified quotient $\mathfrak p_2/2\mathfrak p_2$ is nonreduced, whereas
$\OO_{K_{12}}/\mathfrak P_2\cong\F_4$ is a field.
\end{remark}

\section{The distinguished $n=11$ specialization}

The construction above does not require $n=11$.  It enters afterward as a
distinguished specialization.  Geometrically,
\[
(11,1),(1,11)
\]
lie on $xy=11$ and $x+y=12$, with
\[
(X,Y,T)=(\pm5,\sqrt{11},6).
\]

Arithmetically,
\[
N_{F/\Q}(1+2\sqrt3)=-11.
\]
Let
\[
\mathfrak p_{11}=(1+2\sqrt3).
\]
Since
\[
(1+\sqrt3)(1+2\sqrt3)=7+3\sqrt3
\]
is totally positive,
\[
\boxed{[\mathfrak p_{11}]=[\mathfrak p_2]}
\]
in the narrow ideal class group.  For the narrow Hilbert class field
$K_{12}=\Q(\zeta_{12})$ of $F=\Q(\sqrt3)$,
\[
\boxed{\Art_{K_{12}/F}(\mathfrak p_{11})=T_{11}.}
\]
Here $T_{11}\in\Gal(K_{12}/F)$ because
\[
T_{11}(\sqrt3)=\sqrt3,
\]
and, since $11\equiv-1\pmod{12}$,
\[
T_{11}(\zeta_{12})
=\zeta_{12}^{11}
=\zeta_{12}^{-1}
=\overline{\zeta_{12}}.
\]
Thus $T_{11}$ is complex conjugation over $F$.

In the companion operator representation this involution corresponds to $J$:
\[
\boxed{T_{11}\longleftrightarrow J.}
\]
This is a representation correspondence, not a literal equality between a
field automorphism and an operator in another category.  The specialization
chain is
\[
\boxed{
11\longrightarrow\mathfrak p_{11}
\longrightarrow[\mathfrak p_{11}]=[\mathfrak p_2]
\longrightarrow T_{11}\longleftrightarrow J.}
\]

\section{Symmetry layers and scope}

The finite groups appearing above have different roles:
\[
\begin{array}{c|c}
U(12)\cong V_4 & \text{cyclotomic Galois labels},\\
V_4^{\rm tr} & \text{Boolean translations},\\
V_4^{\rm tr}\rtimes C_2\cong D_8 & \text{translations plus factor exchange},\\
\F_4^\times\cong C_3 & \text{nonzero residue-field multiplication},\\
\operatorname{AGL}(1,4)\cong A_4 & \text{affine field symmetry},\\
\operatorname{A\Gamma L}(1,4)\cong S_4 & \text{affine semilinear symmetry}.
\end{array}
\]
They are compatible layers, not interchangeable groups.

A separate analytic companion may realize the same Pell unit through a
certified operator crossing.  That realization is not used here.  This paper
proves no statement about zeta-zero locations and makes no claim concerning
the Riemann hypothesis.

\section{Conclusion}

Centered Cayley reciprocity selects trace $4$ before arithmetic data enter.
The positive integral normalization then exposes the Pell unit $2+\sqrt3$.
On the ramified ideal it is the integral return $g_{12}$; after centering it
is multiplication by $\sqrt3$; on the real factor cone it is a Lorentz boost;
and after cyclotomic completion it becomes multiplication by $\zeta_{12}$ on
an integral rank-$4$ ideal.  Ramified reduction separates the structure into
order-$3$ multiplication and order-$2$ Frobenius on $\F_4$, while the reduced
Lorentz norm becomes the field trace and hence Boolean XOR.  The $n=11$
shell is then a distinguished example in which the outer $x+y=12$ geometry,
the ramified narrow class, cyclotomic complex conjugation $T_{11}$, and the
companion involution $J$ meet without being used to define the theory.

\end{document}
